\chapter{VOCI DI APOCALISSE GIUNGONO NEL BERKSHIRE}\label{voci-di-apocalisse-giungono-nel-berkshire}

Dopo la festa con gli slittini non vidi Jason per diversi anni pur
tenendomi in contatto con lui. Ci rincontrammo l'anno in cui mi laureai
in medicina, in una casa affittata per l'estate nel Berkshire, a circa
venti minuti da Tanglewood.

Ero stato impegnato. Avevo fatto quattro anni di università più un
periodo di volontariato in una clinica locale e avevo iniziato a
prepararmi per il test di ammissione a medicina un paio di anni prima di
sostenerlo. La media dei voti, i risultati del test e un fascio di
lettere di referenza da parte di tutor universitari e altre personalità
venerabili (oltre alla munificenza di E.D.) mi avevano fatto guadagnare
l'ammissione al campus di medicina dell'Università Statale di New York a
Stony Brook per altri quattro anni. Ormai era terminato, me l'ero
lasciato alle spalle, era finito ma cercavo ancora la possibilità di
svolgere un internato per almeno altri tre anni prima di sentirmi pronto
a esercitare la professione.

Il che mi annoverava fra quelle persone, la maggioranza, che
continuavano a condurre la loro esistenza come se la fine del mondo non
fosse stata annunciata.

Sarebbe stato diverso se avessimo conosciuto l'ora e il giorno esatti
del giorno del giudizio. Ognuno avrebbe potuto scegliere le motivazioni
del proprio personaggio, il panico o la devota rassegnazione, e recitare
fino in fondo la storia dell'umanità con un dignitoso senso del ritmo e
un occhio fisso all'orologio.

Ma ciò che dovevamo affrontare era soltanto la forte probabilità di
un'estinzione finale, in un sistema solare che diventava rapidamente
inadatto alla vita. Forse niente avrebbe potuto proteggerci all'infinito
dal Sole in espansione che tutti avevamo visto nelle immagini della NASA
catturate dalle sonde in orbita\ldots{} eppure, per allora ne eravamo al
riparo per ragioni che nessuno comprendeva. La crisi, se \emph{c'era}
una crisi, era intangibile; l'unica prova concreta era l'assenza delle
stelle - un'assenza come prova, una prova di assenza.

Dunque, come ci si costruisce una vita sotto la minaccia
dell'estinzione? Il quesito determinò la nostra generazione. Per Jason,
a prima vista, fu abbastanza semplice. Si era gettato a capofitto nel
problema: lo Spin stava rapidamente \emph{diventando} la sua vita. E fu
abbastanza facile anche per me, credo. Ero sempre stato votato alla
medicina, e in quell'atmosfera di crisi latente parve una scelta ancora
più saggia. Forse mi ero immaginato a salvare vite nel caso la fine del
mondo si fosse dimostrata più che ipotetica e meno che istantanea. Ma
aveva una qualche importanza se eravamo comunque tutti condannati?
Perché salvare una vita se la vita umana era destinata a essere spazzata
via? Noi dottori in realtà non \emph{salviamo} le vite, noi le
allunghiamo; e se non riusciamo, forniamo cure palliative e sollievo dal
dolore. Il che poteva rivelarsi la capacità più utile di tutte.

Inoltre, l'università e la facoltà di medicina erano state una lunga,
inesorabile e faticosa, ma gradita, distrazione dalle disgrazie del
resto del mondo.

E così io me la cavai. Jason se la cavò. Ma molte persone se la
passarono peggio. Diane fu una di queste.

ooo

Stavo pulendo il monolocale preso in affitto a Stony Brook quando chiamò
Jason.

Era pomeriggio presto. L'illusione ottica indistinguibile dal Sole
splendeva luminosa. Avevo già caricato i bagagli sulla Hyundai, pronto
per tornare a casa. La mia intenzione era di trascorrere un paio di
settimane con mia madre e poi viaggiare per il paese in relax una decina
di giorni. Era l'ultima occasione in cui avrei avuto un po' di tempo
libero prima di iniziare l'internato a Harborview, Seattle e intendevo
usarlo per vedere il mondo, o almeno quella parte racchiusa tra il Maine
e lo stato di Washington. Ma Jason aveva altri programmi. Ebbi giusto il
tempo di dire un ``ciao-come-stai'' che si lanciò nel suo discorso di
convincimento.

«Tyler,» disse, «è un'occasione troppo buona per lasciarsela sfuggire.
E.D. ha affittato una casa estiva nel Berkshire.»

«Davvero? Buon per lui.»

«Ma non può andarci. La scorsa settimana girava per uno stabilimento di
estrusione dell'alluminio in Michigan, è caduto da una piattaforma di
carico e si è rotto l'anca.»

«Mi spiace.»

«Non è una cosa seria, si sta riprendendo ma per un po' deve camminare
con le stampelle e non vuole farsi tutto il viaggio fino al
Massachusetts solo per starsene seduto a oziare e succhiare
antidolorifici. E soprattutto Carol non muore dalla voglia di andarci.»
Non mi sorprendeva per niente. Carol era diventata un'ubriacona
professionista. Non riuscivo a immaginare cosa avrebbe potuto fare nel
Berkshire con E.D. Lawton se non bere un po' di più. «Il fatto è che,»
proseguì Jase, «non può recedere dal contratto, perciò la casa rimarrà
vuota per tre mesi. E quindi pensavo, visto che ti sei laureato, che
potremmo starci insieme per almeno un paio di settimane. Magari potremmo
convincere Diane a unirsi a noi. Vedere un concerto. Camminare nel
bosco. Stare insieme come ai vecchi tempi. In effetti ora sto andando
proprio lì. Che ne dici, Tyler?»

Stavo per rifiutare la sua proposta ma poi pensai a Diane. Mi tornarono
in mente le poche lettere e telefonate che ci eravamo scambiati nelle
occasioni più ovvie e tutte le domande senza risposta accumulate tra di
noi. Sapevo che la cosa più saggia sarebbe stata declinare l'invito ma
era troppo tardi: la bocca aveva già detto di sì.

ooo

E così trascorsi un'altra notte a Long Island; poi stipai le ultime cose
che possedevo al mondo nel bagagliaio dell'auto e seguii la superstrada
turistica del Nord verso l'autostrada di Long Island.

Il traffico era scarso e il tempo incredibilmente piacevole. Era un
pomeriggio dal cielo azzurrissimo, appena caldo. Avrei voluto vendere
l'indomani al miglior offerente e rimanere per sempre al due di luglio.
Mi sentivo scioccamente felice come non lo ero da tempo.

Poi accesi la radio.

Ero abbastanza grande da ricordare i tempi in cui una ``stazione radio''
era un edificio con un trasmettitore e un'antenna a torre, quando la
ricezione radio variava da città a città. Molte di quelle stazioni
esistevano ancora ma la radio analogica della Hyundai era morta circa
una settimana dopo la fine della garanzia, per cui non restava che la
programmazione digitale (ritrasmessa da uno o più aerostati di E.D.
nell'alta atmosfera). In genere ascoltavo brani scaricati di jazz del
ventesimo secolo, una preferenza musicale che avevo acquisito frugando
nella collezione di dischi di mio padre. Mi piaceva pensare che fosse
questa la vera eredità che mi aveva lasciato: Duke Ellington, Billie
Holiday, Miles Davis, musica già vecchia anche ai tempi in cui Marcus
Dupree era giovane, tramandata di nascosto come un segreto di famiglia.
In quel momento desideravo tanto ascoltare qualche pezzo jazz, ma il
tizio che mi aveva revisionato la macchina prima del viaggio aveva
cancellato tutte le stazioni memorizzate e si era sintonizzato su un
canale che non riuscivo a cambiare. Perciò ero costretto a sentire di
disastri naturali e pettegolezzi sulle celebrità. Si parlava anche dello
Spin.

Era stato in quel periodo che avevamo iniziato a chiamarlo ``Spin''.

Anche se la maggior parte del mondo non credeva che esistesse.

I sondaggi parlavano piuttosto chiaro. La NASA aveva reso pubblici i
dati delle sonde orbitali la sera in cui Jason aveva svelato la notizia
a me e Diane, e i lanci europei avevano confermato i risultati
americani. Eppure, anche otto anni dopo che lo Spin era diventato di
pubblico dominio, solo una minoranza di europei e nordamericani lo
considerava ``una minaccia per loro o per le loro famiglie''. In molte
zone dell'Asia, Africa e Medio Oriente, robuste maggioranze
consideravano il tutto un complotto o un incidente degli Stati Uniti,
probabilmente un tentativo fallito di creare un qualche genere di scudo
spaziale.

Una volta ne avevo chiesto la ragione a Jason. Mi aveva risposto: «Pensa
a cosa stiamo chiedendo di credere. Parliamo, a livello globale, di una
popolazione con una conoscenza prenewtoniana dell'astronomia. È davvero
necessario conoscere certe cose della Luna e delle stelle se passi la
vita a rimediare biomassa sufficiente a sfamare te e la tua famiglia?
Per raccontare a quelle persone qualcosa di significativo sullo Spin
dovresti partire da molto lontano. Tanto per cominciare, dovresti dire
che la Terra ha qualche miliardo di anni. Si dovrebbero misurare, forse
per la prima volta, con il concetto di ``miliardo di anni''. È tanta
roba da digerire, soprattutto se sei cresciuto in una teocrazia
mussulmana, un villaggio animista o una scuola pubblica della
\emph{Bible Belt}. Poi dovresti spiegare che la Terra non è immutabile,
che c'è stata un'era più lunga della nostra in cui gli oceani erano
vapore e l'aria veleno. Che le forme di vita si sono sviluppate
spontaneamente e si sono evolute casualmente per tre miliardi di anni
prima di dare vita a un essere che si possa definire umano. Poi dovresti
parlare del Sole, del fatto che neppure il Sole è stabile ma che
all'inizio era una nube di gas e polvere addensati e che un giorno, fra
qualche altro miliardo di anni da ora, si espanderà, inghiottirà la
Terra e alla fine espellerà i suoi strati esterni e si ridurrà a una
pallina di materia superdensa. Le basi della cosmologia, no? Tu le hai
apprese da tutti quei libri che leggevi, per te è una seconda pelle, ma
per la maggior parte della gente è una visione del mondo del tutto nuova
e forse offensiva per molti dei loro principi. Perciò devi aspettare che
queste cose facciano presa. Una volta che hanno fatto presa, poi devi
dare la \emph{vera} brutta notizia: il \emph{tempo stesso} è fluido e
imprevedibile. Il mondo che sembra così normale, nonostante quello che
abbiamo appena imparato, di recente è stato chiuso in una specie di
cella frigorifera cosmologica. \emph{Perché} ci hanno fatto questo? Non
lo sappiamo di preciso. Crediamo che provenga dall'azione deliberata di
entità potenti e inaccessibili, tanto da poterli chiamare dei. E se
facciamo adirare gli dei, potrebbero toglierci la loro protezione, e ben
presto le montagne fonderebbero e gli oceani bollirebbero. Ma non è che
dobbiate crederci sulla parola. Lasciate perdere i tramonti e la neve
che come ogni inverno arriva in montagna. Abbiamo le prove. Abbiamo
calcoli, deduzioni logiche e scatti fotografici presi da videocamere.
Prove forensi di altissimo valore.» Jason aveva mostrato uno dei suoi
sorrisi tristi e beffardi. «Strano\ldots{} la giuria non è convinta.»

E non erano solo gli ignoranti a non essere convinti. Alla radio, un
amministratore delegato del settore assicurativo aveva cominciato a
lamentarsi dell'impatto economico di ``tutta questa discussione
inflessibile, acritica sul cosiddetto Spin,'' per poi sentenziare, ``la
gente comincia a crederci davvero.'' E questo era un male per
l'economia. Rendeva le persone imprudenti. Incoraggiava l'immoralità, il
crimine e il deficit. Ancor peggio, incasinava le tavole attuariali.
``Se il mondo non finirà nei prossimi trenta o quarant'anni, potremmo
andare incontro a un disastro.''

Le nuvole cominciarono a addensarsi da ovest. Un'ora dopo quel bel cielo
azzurro si era annuvolato in modo uniforme e la pioggia cominciò a
battere sul parabrezza. Accesi i fari.

Il notiziario alla radio passò oltre. Discussero a lungo di un'altra
notizia riportata dai titoli recenti dei giornali: le casse argentee.
Erano grandi quanto città e fluttuavano oltre la barriera dello Spin, a
centinaia di chilometri sopra i poli terrestri. Fluttuavano, non
orbitavano. Un oggetto sospeso in un'orbita stabile sull'equatore - come
i satelliti geosincroni - sarebbe stata una cosa comprensibile, ma
secondo le più elementari leggi del moto niente dovrebbe ``orbitare'' in
una posizione fissa sul polo del pianeta. Eppure quegli oggetti si
trovavano lì, rilevati da una sonda radar e fotografati di recente da
una missione automatizzata di passaggio ravvicinato: un altro strato del
mistero dello Spin, altrettanto incomprensibile per le masse incolte, in
questo caso specifico anche per me. Avrei voluto parlarne con Jason.
Forse desideravo che fosse lui a dare un senso a tutto questo al posto
mio.

ooo

Pioveva a dirotto, i tuoni rimbombavano fra le colline quando finalmente
mi fermai vicino alla casa che E.D. Lawton aveva preso in affitto fuori
Stockbridge per un breve periodo.

La proprietà era un cottage inglese in stile rustico con quattro stanze
e le mura esterne pitturate di verde arsenico, situata all'interno di un
bosco protetto di una quarantina di ettari. Brillava al crepuscolo come
una lampada antivento. Jason era già lì, la sua Ferrari bianca era
parcheggiata sotto una tettoia che andava dalla casa al garage.

Doveva avermi sentito arrivare: aprì la grande porta d'ingresso prima
che bussassi. «Tyler!» esclamò sorridendo.

Entrai e poggiai la mia unica valigia inumidita di pioggia sul pavimento
piastrellato dell'ingresso. «Ne è passato di tempo» dissi.

Eravamo rimasti in contatto per e-mail e telefono ma salvo un paio di
brevi comparse durante le vacanze alla Grande Casa, questa era la prima
volta che ci trovavamo nella stessa stanza in quasi otto anni. Suppongo
che il tempo passato mostrasse su entrambi un lieve repertorio di
cambiamenti. Avevo dimenticato quanto fosse formidabile d'aspetto. Era
sempre stato alto, sempre disinvolto nel suo corpo; lo era ancora, anche
se sembrava più magro, non delicato ma delicatamente bilanciato, come
una scopa in equilibrio sulla punta del manico. Aveva i capelli lunghi
circa mezzo centimetro e ispidi come stoppie. E malgrado guidasse una
Ferrari, era completamente a digiuno di stile personale: indossava jeans
strappati, un maglione abbondante fatto ai ferri e punteggiato di
palline sfilacciate, e scarpe da ginnastica economiche.

«Hai mangiato per strada?» mi chiese.

«Ho pranzato tardi.»

«Hai fame?»

Non ne avevo ma ammisi di desiderare tanto una tazza di caffè. La
facoltà di medicina mi aveva fatto diventare un caffeinomane. «Sei
fortunato,» rispose Jason, «lungo la strada ho comprato mezzo chilo di
guatemalteco.» I guatemaltechi, indifferenti alla fine del mondo,
coltivavano ancora caffè. «Metto su la caffettiera. Nel frattempo ti
mostro la casa.»

Facemmo un giro. Aveva un'aria stipata tipica del ventesimo secolo,
pareti di colore verde mela o arancia acerba, mobili antichi solidi e
rustici e letti in ottone, tende di pizzo sopra vetri bombati su cui la
pioggia scorreva ininterrottamente. Comfort moderni in cucina e in
soggiorno, un grande televisore, un impianto stereo, una connessione
internet. Accogliente nella pioggia. Ritornati al piano di sotto, Jason
servì il caffè. Ci sedemmo al tavolo della cucina e cercammo di
recuperare il tempo perduto.

Riguardo al lavoro Jase rimase sul vago, per modestia o ragioni di
sicurezza. Negli otto anni trascorsi dalla rivelazione della vera natura
dello Spin era riuscito a ottenere un dottorato in astrofisica ma poi lo
aveva abbandonato per ricoprire un incarico di primo livello nella
Fondazione Perielio di E.D. Forse non era stata una cattiva mossa, visto
che E.D. era un membro importante della Commissione d'inchiesta del
Presidente Walker sulla Pianificazione della crisi ambientale e globale.
Secondo Jase, la Perielio sarebbe passata da gruppo di esperti
aerospaziali a organo consultivo ufficiale con una vera e propria
autorità nel determinare la linea politica.

Chiesi: «È legale?»

«Non essere ingenuo, Tyler. E.D. si è già staccato dalla Compagnia
Lawton. Ha dato le dimissioni dal Consiglio di amministrazione e le sue
azioni sono gestite da un blind trust. Secondo i nostri avvocati non c'è
conflitto d'interessi.»

«E quindi tu che fai alla Perielio?»

Sorrise. «Ascolto con attenzione i miei superiori,» disse, «e faccio
proposte garbate. Raccontami della facoltà di medicina.»

Mi chiese se avessi provato disgusto a vedere tanta debolezza e malattia
umana. E così gli raccontai del corso di anatomia al secondo anno. Io e
un'altra decina di studenti avevamo dissezionato un cadavere e
classificato gli organi interni per dimensione, colore, funzione e peso.
Non era stata affatto un'esperienza piacevole. L'unica consolazione era
stata la sua autenticità e l'unica virtù l'utilità. Ma era stato anche
un segnale di passaggio. Oltre quel punto, non era rimasto più nulla
dell'infanzia.

«Cavolo, Tyler. Vuoi qualcosa di più forte del caffè?»

«Non sto dicendo che è stato un dramma. È questa la cosa scioccante.
\emph{Non è stato un} dramma. Esci da lì e te ne vai a vedere un film.»

«Ne hai fatta di strada dalla Grande Casa, comunque.»

«Tanta. Ne abbiamo fatta entrambi.» Sollevai la tazza.

Poi ci abbandonammo ai ricordi e la tensione svanì dai nostri discorsi.
Parlammo dei vecchi tempi. Ricademmo nel solito schema. Jason menzionava
un luogo - il seminterrato, il centro commerciale, il ruscello nel bosco
- e io ci infilavo la storia: quando forzammo l'armadietto dei liquori;
quando vedemmo una ragazza della Rice di nome Kelley Weems rubare una
scatola di preservativi in una farmacia; l'estate in cui Diane
insistette per leggerci dei passaggi angosciosi di Christina Rossetti,
come se avesse scoperto qualcosa di profondo.

«Il grande prato» offrì Jason.

«La notte in cui scomparvero le stelle» rilanciai io.

E poi rimanemmo in silenzio per un po'.

Alla fine dissi: «Allora\ldots{} viene o no?»

«Ci sta pensando» rispose Jase in tono neutro. «Si sta destreggiando fra
gli impegni. Dovrebbe chiamare domani e farmi sapere.»

«È ancora giù al Sud?» Era la notizia più recente che avevo sentito,
riferita da mia madre. Diane era in una qualche università del Sud a
studiare qualcosa che non riuscivo proprio a ricordare: geografia
urbana, oceanografia o qualche altra improbabile -ografia.

«Sì, ancora lì» confermò Jason, sistemandosi sulla sedia. «Sai, Ty, sono
cambiate tante cose per Diane.»

«Non credo ci sia da meravigliarsi.»

«È mezza fidanzata. Sta per sposarsi.»

Reagii con un certo decoro. «Be', buon per lei.» Come avrei mai potuto
essere geloso? Non avevo più rapporti con Diane - non ne avevo mai avuto
uno in quel senso particolare della parola ``rapporto''. Io stesso a
Stony Brook mi ero quasi fidanzato con una studentessa del secondo anno
di nome Candice Boone. Ci eravamo gustati i ``ti amo'' detti, almeno
fino a quando la cosa non ci aveva stancato. Mi sembra che si fosse
stufata per prima Candice.

E comunque: \emph{mezza} fidanzata? Come funzionava questa cosa?

Ero tentato di chiederglielo. Ma Jason era chiaramente a disagio per la
piega che aveva preso la conversazione, il che mi risvegliò un ricordo:
una volta, ai tempi della Grande Casa, Jason aveva portato una ragazza
per farle conoscere la sua famiglia. Era insignificante ma gradevole,
l'aveva conosciuta al club di scacchi della Rice, ed era molto timida,
parlava poco. Quella sera Carol era rimasta abbastanza sobria, ma la
disapprovazione di E.D. era stata evidente: si era comportato in maniera
molto rude e quando la ragazza se n'era andata aveva rimproverato Jase
per aver ``trascinato in casa un esemplare del genere''.

``A un grande intelletto si deve accompagnare una grande
responsabilità'' aveva detto E.D., che non voleva che Jason si
ritrovasse incastrato con l'inganno in un matrimonio ordinario. Non
voleva vederlo finire a ``cambiare pannolini'' quando invece poteva
``lasciare un segno nel mondo''.

Al posto di Jason in molti avrebbero smesso di portare a casa le persone
che frequentavano.

Jason aveva smesso del tutto di frequentare ragazze.

ooo

La casa era vuota quando mi svegliai il mattino dopo.

Trovai un appunto sul tavolo della cucina: ``Torno a mezzogiorno o più
tardi.'' Jase era uscito a prendere qualcosa per una grigliata. Erano le
nove e mezzo. Avevo dormito magnificamente e fino a tardi, pervaso del
classico languore da vacanze estive.

Sembrava fosse la casa a generarlo. I temporali della sera precedente
erano passati e una gradevole brezza mattutina filtrava dalle tende in
calicò. La luce del giorno metteva in risalto le imperfezioni nelle
venature dei ripiani in legno della cucina. Feci colazione con calma
vicino alla finestra guardando le nuvole che come maestose golette
veleggiavano nell'orizzonte.

Poco dopo le dieci suonarono al campanello e per un attimo mi prese il
panico al pensiero che potesse essere Diane: aveva deciso di presentarsi
in anticipo? Ma poi scoprii che era Mike-il-giardiniere in bandana e
maglietta smanicata, venuto ad avvisare che stava andando a sistemare il
prato. Non voleva svegliare nessuno, ma il tosaerba era piuttosto
rumoroso. Poteva tornare nel pomeriggio se era un problema. «Nessun
problema» dissi, e qualche minuto dopo percorreva il perimetro della
proprietà su un vecchio John Deere verde che affumicava l'aria con il
puzzo di olio bruciato. Ancora un po' assonnato, mi chiesi come sarebbe
parso questo lavoro di giardinaggio agli occhi di quello che Jason amava
chiamare l'Universo in generale. Per l'Universo al di là dello Spin, la
Terra era un pianeta quasi in stasi. Quei fili d'erba erano cresciuti
attraverso i secoli, in un movimento maestoso quanto l'evoluzione delle
stelle. Mike, una forza della natura nata un paio di miliardi di anni
prima, li falciava con grande e irresistibile pazienza. I fili recisi
cadevano come sfiorati lievemente dalla gravità, mentre equinozi e
solstizi passavano incessanti e vermi Matusalemme strisciavano tra la
terra e il concime mentre altrove nella galassia potevano sorgere e
cadere imperi.

Naturalmente Jason aveva ragione: era una cosa difficile da credere.
Anzi, no, non ``da credere'' - la gente crede a qualsiasi cosa di non
plausibile - ma da accettare come verità fondamentale sul mondo. Mi
sedetti sotto il portico della casa, distante quanto più possibile dal
trattore rumoroso. L'aria era fredda e il Sole piacevole sul viso quando
mi girai verso di esso, anche se era solo la radiazione filtrata dallo
Spin di una stella che andava a tutta velocità fuori dalla sua membrana,
in un mondo in cui i secoli venivano sperperati come secondi.

Non può essere vero. \emph{È} vero.

Ripensai alla facoltà di medicina, al corso di anatomia di cui avevo
parlato a Jason. Candice Boone, la mia ex quasi fidanzata, aveva
frequentato quel corso con me. Era stata stoica durante la dissezione ma
non dopo. «Un corpo umano» aveva detto «dovrebbe contenere amore, odio,
coraggio, viltà, anima, spirito\ldots{} non questo assortimento viscido
di imponderabili azzurri e rossi.» Sì. E non dovremmo essere trascinati
contro la nostra volontà in un futuro duro e mortale.

«Ma il mondo è ciò che è e non ci si può contrattare.» Questo è quanto
le avevo risposto.

Mi disse che ero freddo. Eppure erano ancora le parole più sagge che
fossi mai riuscito a mettere insieme.

ooo

La mattinata volò via. Mike finì di sistemare il prato e se ne andò
lasciando l'aria piena di silenzio umido. Dopo un po' mi diedi una mossa
e telefonai a mia madre in Virginia, dove il tempo, mi raccontò, era
meno allettante che nel Massachusetts: era ancora nuvoloso dopo la
tempesta della notte precedente che aveva abbattuto alcuni alberi e
qualche linea elettrica. Le dissi che ero arrivato sano e salvo alla
casa estiva di E.D. Mi chiese come avessi trovato Jason, anche se lei
forse lo aveva visto più di recente, durante una delle sue visite alla
Grande Casa. «Più vecchio ma lo stesso Jase di sempre» le risposi.

«È preoccupato per quest'affare della Cina?»

Dall'Evento di ottobre mia madre era diventata una maniaca delle ultime
notizie, guardava la CNN non per piacere e nemmeno per tenersi
informata: lo faceva soprattutto per rassicurarsi, nello stesso modo in
cui un abitante di un villaggio messicano tiene d'occhio un vulcano
vicino, sperando di non vederlo fumare.

«In questa fase la questione cinese è solo una crisi diplomatica,» mi
spiegò, «anche se una velata minaccia c'è stata. Qualcosa che ha a che
fare con il progetto controverso di un lancio satellitare. Dovresti
chiederlo a Jason.»

«È E.D. che ti tormenta con questa roba?»

«Macché! Ogni tanto sento dire qualcosa da Carol.»

«Non so quanto ti possa fidare.»

«Ma va', Ty. Beve ma non è mica stupida. E neanch'io, per la
precisione.»

«Non intendevo dire quello.»

«Quasi tutto quello che sento in questi giorni su Jason e Diane proviene
dalla bocca di Carol.»

«Ha detto se Diane sta arrivando nel Berkshire? Da Jase non riesco mai
ad avere risposte dirette.»

Mia madre esitò. «Negli ultimi due anni Diane è stata un po'
imprevedibile. Immagino che si tratti solo di questo.»

«In che senso ``imprevedibile''?»

«Be', sai, a scuola non va tanto bene. Un po' di problemi con la
giustizia\ldots»

«Con la \emph{giustizia}?»

«No, non è che abbia rapinato una banca o roba del genere ma è stata
arrestata un paio di volte quando le manifestazioni dell'NR sono
sfuggite di mano.»

«E che cavolo ci faceva alle manifestazioni dell'NR?»

Un'altra pausa. «Dovresti proprio chiederlo a Jason.»

Intendevo farlo.

La sentii tossire - me la immaginai con la mano sul telefono, la testa
girata delicatamente - e chiesi: «Come ti senti?»

«Stanca.»

«Qualche novità dal dottore?» Stava facendo delle cure per l'anemia.
Flaconi di compresse di ferro.

«No. Sto solo invecchiando, Ty. Tutti invecchiamo prima o poi.» Quindi
aggiunse: «Sto pensando di andare in pensione. Se quello che faccio puoi
chiamarlo lavoro. Ora che i gemelli se ne sono andati, ci sono solo
Carol e E.D., e lui neanche tanto, da quando è iniziata la faccenda di
Washington.»

«A loro lo hai detto?»

«Non ancora.»

«Non sarebbe più la Grande Casa senza di te.»

Rise ma non di felicità. «Penso che averle dedicato una vita intera sia
più che sufficiente, grazie.»

Ma non fece più cenno al fatto di andarsene. Credo che fu Carol a
convincerla a restare.

ooo

Jase sbucò dalla porta d'ingresso a metà pomeriggio. «Ty?» I jeans
enormi gli si appendevano ai fianchi come il cordame di una nave in
bonaccia e la maglietta era schizzata di fantasmi di macchie di salsa.
«Puoi darmi una mano con il barbecue?»

Andai con lui sul retro della casa. Il barbecue era una normale griglia
a gas propano. Jase non ne aveva mai usata una. Aprì la valvola del
serbatoio, schiacciò il pulsante di accensione ed ebbe un sussulto
quando la fiamma si animò. Poi sogghignando mi disse: «Abbiamo delle
bistecche. E un'insalata ai tre fagioli presa nella gastronomia in
città.»

«E quasi nessuna zanzara.»

«Hanno fatto la disinfestazione questa primavera. Hai fame?»

Ne avevo. In qualche modo, sonnecchiando nel pomeriggio, mi si era
stimolato l'appetito. «Cuciniamo per due o per tre?»

«Sto ancora aspettando che Diane mi chiami. Lo farà stasera, forse.
Credo che a cena saremo solo io e te.»

«Sempre che i cinesi non ci facciano arrosto prima.»

Era un'esca.

Jason abboccò. «Sei preoccupato per i \emph{cinesi}, Ty? Non c'è neanche
più la crisi. È stata risolta.»

«Che sollievo.» Nella stessa giornata ero venuto a conoscenza della
crisi e della sua pronta risoluzione. «Me l'ha accennato mia madre.
Qualcosa che ha sentito al notiziario.»

«L'esercito cinese vuole distruggere i manufatti sul polo. Hanno dei
missili a testata nucleare sulle piattaforme di Jiuquan, pronti per il
lancio. Il ragionamento è questo: se riescono a danneggiare i
dispositivi sul polo potrebbero abbattere anche l'intero scudo di
ottobre. Naturalmente non c'è ragione di credere che funzioni. Che
probabilità abbiamo che una tecnologia capace di manipolare il tempo e
la gravitazione sia vulnerabile alle nostre armi?»

«Quindi abbiamo minacciato i cinesi e loro hanno fatto marcia indietro?»

«Una cosa del genere. Ma gli abbiamo fatto anche un'offerta allettante.
Gli abbiamo offerto di salire a bordo con noi.»

«Non capisco.»

«Gli diamo la possibilità di unirsi a noi nel piccolo progetto di
salvare il mondo.»

«Ora mi spaventi un po', Jase.»

«Passami quelle pinze. Mi spiace, so che può sembrare criptico. Non
dovrei proprio parlare di queste cose. Con nessuno.»

«Stai facendo un'eccezione per me?»

«Per te faccio sempre un'eccezione.» E sorridendo aggiunse: «Ne
discutiamo a pranzo, d'accordo?»

Lo lasciai vicino alla griglia avvolto nel fumo e nel calore.

ooo

La stampa aveva accusato due amministrazioni americane consecutive di
non avere fatto nulla per lo Spin. Ma era una critica priva di
fondamento. Seppure ci fosse stata la \emph{possibilità} di fare
qualcosa di concreto, nessuno sapeva quale fosse. E qualsiasi azione di
evidente rappresaglia - come quella proposta dai cinesi - sarebbe stata
molto pericolosa.

La Perielio spingeva per un approccio radicalmente diverso.

«La metafora più appropriata,» disse Jason, «non è il combattimento ma
il judo. Sfruttare il peso e la velocità di un avversario più grosso a
suo discapito. È questo che vogliamo fare con lo Spin.»

Me lo disse in maniera concisa mentre tagliava la bistecca grigliata con
attenzione chirurgica. Mangiammo in cucina con la porta sul retro
aperta. Un bombo enorme, talmente grasso e giallo da sembrare un nodo di
fili di lana trasportato dall'aria, andò a sbattere contro la
zanzariera.

«Prova a pensare allo Spin come a un'opportunità anziché a un attacco.»

«Un'opportunità per fare cosa? Morire prematuramente?»

«Un'opportunità di usare il tempo per i nostri scopi, in un modo in cui
finora non abbiamo potuto.»

«Non è proprio il tempo quello che ci hanno portato via?»

«Al contrario. Fuori dalla nostra bollicina terrestre abbiamo milioni di
anni con cui giocare. E abbiamo uno strumento che funziona in maniera
estremamente affidabile proprio su quelle grandezze temporali.»

«Strumento» brontolai perplesso, mentre lui infilzava un altro
quadratino di manzo. Il pranzo era frugale. Una bistecca su un piatto,
una bottiglia di birra accanto. Senza fronzoli, tranne che per
l'insalata ai tre fagioli di cui prese una porzione modesta.

«Sì, uno strumento, quello più ovvio: l'evoluzione.»

«L'evoluzione.»

«Tyler, non possiamo discutere se ti limiti a ripetermi tutto quello che
ti dico.»

«D'accordo, bene, l'evoluzione è uno strumento\ldots{} però non riesco
ancora a capire come in trenta o quarant'anni possiamo evolverci in modo
sufficiente da cambiare le cose.»

«Non noi, \emph{diamine}, e certo non in trenta o quarant'anni. Parlo di
forme di vita semplici. Parlo di eoni. Parlo di Marte.»

«Marte.» Ops.

«Non essere ottuso. Pensaci.»

Marte era un pianeta morto dal punto di vista funzionale, anche se un
tempo poteva avere avuto precursori primitivi di forme di vita. Al di
fuori della bolla dello Spin, dall'Evento di ottobre, Marte si era
``evoluto'' per milioni di anni, riscaldato dal Sole in espansione.
Secondo le ultime fotografie orbitali, però, rimaneva un pianeta arido,
morto. Se avesse posseduto forme di vita semplici e un clima favorevole,
a quell'ora sarebbe dovuto diventare una giungla verde e rigogliosa. Ma
non era accaduto niente di tutto questo.

«La gente parlava di terraformazione» disse Jason. «Ti ricordi quei
romanzi di fantascienza che leggevi?»

«Li leggo ancora, Jase.»

«Buon per te. Come faresti a terraformare Marte?»

«Cercherei di introdurre sufficiente gas serra nell'atmosfera per
riscaldarla. Di far sciogliere la sua acqua ghiacciata. Ci pianterei
degli organismi semplici. Ma anche con le ipotesi più ottimistiche, ci
vorrebbe\ldots»

Lui sorrise.

«Mi stai prendendo in giro» mi lamentai.

«No.» Il sorriso svanì. «Niente affatto. No, è una cosa serissima.»

«Ma da cosa partireste\ldots?»

«Partiremmo da una serie di lanci sincronizzati con carichi di batteri
geneticamente modificati. Semplici motori a ioni e un lento scivolamento
verso Marte. Impatti per lo più controllati a cui possono sopravvivere
gli organismi unicellulari, e qualche carico più grande con testate a
penetrazione profonda per liberare quegli stessi organismi sotto la
superficie del pianeta, dove sospettiamo che ci sia acqua interrata.
Riduciamo i rischi con lanci multipli e un ampio spettro di organismi
candidati. L'idea è di avere sufficiente azione organica per liberare il
carbonio imprigionato nella crosta e disseminarlo nell'atmosfera. Gli
diamo qualche milione di anni - qualche mese per il nostro tempo - e poi
riosserviamo il pianeta. Se è diventato un luogo più caldo con
un'atmosfera più densa e magari qualche pozza di acqua semighiacciata,
ricominciamo il ciclo, stavolta con piante multicellulari modificate
appositamente per quel tipo di ambiente. Queste immetteranno
dell'ossigeno nell'aria e forse la pressione atmosferica aumenterà di un
altro paio di millibar. Ripetere se necessario. Aggiungere altri milioni
di anni e mescolare. In un tempo ragionevole - il tempo misurato dai
nostri orologi, intendo - potremmo riuscire a cucinare un pianeta
abitabile.»

Era un'idea sbalorditiva. Mi sentivo come uno di quei personaggi spalla
di un romanzo giallo vittoriano, uno di quelli che avrebbe detto
qualcosa del genere: «Era un piano audace, persino ridicolo, quello che
aveva congegnato ma, per quanti sforzi avessi fatto, non vi avevo
trovato alcun difetto!»

Eccetto uno. Un difetto fondamentale.

«Jason, ammettiamo che sia possibile, ma a cosa ci \emph{servirebbe}?»

«Se Marte è abitabile, la gente può andarci a vivere.»

«La gente? Tutti i sette o otto miliardi di persone che abitano la
Terra?»

Sbuffò. «Macché, no, solo qualche pioniere. Riproduttori, se vogliamo
metterla in termini scientifici.»

«E cosa dovrebbero fare?»

«Vivere, riprodursi e morire. Milioni di generazioni per ogni anno
terrestre.»

«E a che scopo?»

«Se non altro a dare alla specie umana una seconda possibilità nel
sistema solare. Nel migliore dei casi avranno tutta la conoscenza che
possiamo fornire loro più qualche milione di anni per migliorarla.
All'interno della bolla dello Spin non abbiamo tempo sufficiente per
scoprire chi sono gli Ipotetici o il perché ci stanno facendo questo. I
nostri eredi marziani potrebbero avere una possibilità migliore. Forse
potranno darci loro una risposta.»

O combattere al posto nostro?

(Per inciso, questa era la prima volta che li sentivo chiamare
``Ipotetici'' - le ipotetiche intelligenze controllanti, le creature
invisibili e assai teoriche che ci avevano rinchiuso nella loro cripta
temporale. Quel nome non si sarebbe diffuso tra il grande pubblico che
qualche anno più tardi. Ne fui dispiaciuto quando accadde. La parola era
troppo asettica, evocava qualcosa di astratto e freddamente oggettivo;
probabilmente la verità era più complessa.)

«Ma c'è un progetto per \emph{fare} queste cose concretamente?»

«Oh, sì.» Jason aveva finito tre quarti della sua bistecca. Scostò il
piatto. «Non è neanche troppo costoso. Modificare organismi unicellulari
estremamente resistenti è l'unica parte problematica. La superficie di
Marte è fredda, arida, di fatto senz'aria e investita dalla radiazione
sterilizzante ogni volta che il Sole sorge. Comunque, abbiamo tantissimi
estremofili con cui lavorare - batteri che vivono nelle rocce antartiche
o negli scarichi dei reattori nucleari. Tutto il resto è tecnologia
ampiamente dimostrata. Sappiamo che i missili funzionano. Sappiamo che
l'evoluzione organica funziona. L'unica cosa davvero nuova è la nostra
prospettiva. Essere in grado di ottenere risultati a lunghissimo termine
soltanto a giorni o mesi di distanza dopo il lancio. Si tratta
di\ldots{} la gente la chiama ``ingegneria teleologica''.»

«È simile a quello,» azzardai collaudando la parola che mi aveva donato,
«che stanno facendo gli Ipotetici.»

«Sì» disse Jason sollevando il sopracciglio in un'espressione che dopo
tutti quegli anni trovavo ancora lusinghiera: sorpresa, rispetto. «Sì,
in un certo senso credo che lo sia.»

ooo

Una volta avevo letto un particolare interessante in un libro sul primo
allunaggio con uomini a bordo del 1969. A quel tempo, diceva il libro,
alcuni tra i più anziani - uomini e donne nati nel diciannovesimo
secolo, abbastanza vecchi da ricordare un mondo precedente alle
automobili e alla televisione - erano stati restii a credere alla
notizia. Parole che avevano avuto un senso solo nelle fiabe della loro
infanzia (``Stanotte due uomini hanno camminato sulla Luna'') venivano
date come dichiarazioni di fatto. E loro non riuscivano ad accettarlo.
Confondevano il senso di ciò che era razionale e ciò che era assurdo.

Ora toccava a me.

«Andremo a terraformare e colonizzare Marte» aveva detto il mio amico
Jason, e non stava delirando\ldots{} o perlomeno non era più folle delle
decine di persone intelligenti e potenti che evidentemente condividevano
la sua convinzione. Perciò l'affermazione era seria; almeno a livello
burocratico doveva già essere in corso d'opera.

Dopo pranzo feci una passeggiata per i campi mentre c'era ancora un po'
di luce.

Mike-il-giardiniere aveva fatto un lavoro discreto. Il prato brillava
come l'idea di giardino di un matematico, la coltivazione di un colore
primario. Oltre il prato, le ombre avevano cominciato a crescere negli
ettari boscosi. Diane avrebbe apprezzato il bosco sotto quella luce. Mi
tornarono di nuovo in mente quelle riunioni estive vicino al ruscello,
ormai anni prima, quando ci leggeva passi di libri antichi. Una volta,
parlando dello Spin, Diane aveva citato una piccola poesia del poeta
inglese A.E. Housman:

``\emph{L'orso grizzly, gran gigante},

\emph{ha sbranato or or l'infante}.

\emph{All'infante non è occorso}

\emph{che se l'è mangiato l'orso}.''

ooo

Jason era al telefono, quando tornai passando dalla porta della cucina.
Mi guardò e, voltandomi le spalle, abbassò la voce.

«No, se dev'essere così, ma\ldots{} no, capisco. Va bene. Ho detto ``va
bene'', no? Va bene significa va bene.»

Infilò il telefono in tasca. «Era Diane?» gli chiesi.

Annuì.

«Sta arrivando?»

«Sta arrivando. Ma ci sono un paio di cose di cui ti voglio parlare
prima che arrivi. Hai presente le cose di cui abbiamo parlato durante il
pranzo? Non possiamo condividerle con lei. E con nessun altro, in
realtà. Non è un'informazione pubblica.»

«Vuoi dire che è segretata.»

«Tecnicamente sì, suppongo che sia così.»

«Ma tu \emph{me} ne hai parlato.»

«Sì, è stato un crimine federale.» Aggiunse sorridendo: «Mio, non tuo. E
conto sul fatto che tu sia discreto. Abbi pazienza, fra un paio di mesi
sarà tutto trasmesso dalla CNN. E poi ho dei progetti per te, Ty. Uno di
questi giorni la Perielio esaminerà dei candidati per condurre una vita
rurale molto dura. Avremo bisogno di ogni genere di medico sul posto.
Non sarebbe grandioso se potessi essere uno di questi, se potessimo
lavorare insieme?»

Ero sbigottito. «Mi sono appena laureato, Jase. Non ho ancora fatto
l'internato.»

«Ogni cosa a suo tempo.»

«Non ti fidi di Diane?» gli chiesi.

Il sorriso gli si sgonfiò. «Sinceramente no, non più. Non ultimamente.»

«Quando arriva?»

«Domani, prima di mezzogiorno.»

«E cos'è che non vuoi dirmi?»

«Porta il fidanzato.»

«Ed è un problema?»

«Vedrai.»